
\subsection{Model}
We adopt the flow-matching~\citep{flow-matching} based generative models and the diffusion transformer (DiT)~\citep{dit} model architecture.
Additional conditioning modules are added to provide control.
\cref{fig:audio_model} illustrates the model architecture.

\begin{figure}[h]
    \centering
    \includegraphics[width=\linewidth]{figures/audio/MovieGen_Audio_Model_v4.pdf}
    \caption{\textbf{\mvga{} model diagram.} Yellow blocks denote input, blue blocks denote pre-trained and frozen modules, gray blocks denote operations without learnable parameters, green blocks denote learnable modules, and the pink block shows the output velocity $u(\bx_t, \bc, t; \theta)$. Conditioning input $\bc$ includes masked audio context, video, and text. $\bx_t$ is a sample from $p_t$, and $t$ is the flow time step. For audio context, we replace the masked frames with zeros for DAC-VAE output.
    }
    \label{fig:audio_model}
\end{figure}

\subsubsection{Flow-matching}
We also use the Flow Matching~\citep{flow-matching} objective as described in~\cref{sec:t2v_training_objective} to train the \mvga{} model.
The same optimal transport path is used for constructing $\bx_{t}$ which is now an audio sample in the latent space that we will describe in \cref{sec:audio_repr_and_cond}, and the same logit-normal distribution is used for sampling flow-step $t$.
Instead of conditioning only on text prompt embedding $\bp$, \mvga{} is conditioned on multimodal prompt $\bc$ which will be described in \cref{sec:audio_repr_and_cond}.

We choose diffusion-style models~\citep{Ho2020DenoisingDP, song2020denoising} over discrete token-based language models~\citep{kreuk2022audiogen} because
(1) it shows strong empirical performance in sound, music, and speech generation~\citep{liu2023audioldm, ghosal2023text, majumder2024tango, shen2023naturalspeech, huang2023noise2music},
(2) its non-autoregressive nature permits flexible generation direction, and can be used for both infilling or out-filling in both directions,
(3) modeling audio in continuous space enables applications of techniques such as SDEdit~\citep{meng2021sdedit} for editing and multi-diffusion~\citep{bar2023multidiffusion} for infinite-length audio generation,
(4) it enables users to flexibly trade-off quality for runtime through configuring ODE parameters,
and enjoys recent advancement in distillation or consistency training techniques that boost quality significantly at a much lower runtime.
On the other hand, we choose flow-matching over diffusion because we found it achieves better training efficiency, inference efficiency, and performance compared to diffusion models as shown in recent works~\citep{lan2024high, le2023voicebox, vyas2023audiobox, prajwalmusicflow, mehta2024matcha, sd3}.

\subsubsection{Diffusion Transformer}
\mvga{} adopts the diffusion transformer (DiT) architecture~\citep{dit},
which modulates the  outputs of normalization layers with scale and bias,
and outputs of self-attention and feed-forward layers with scale in each transformer block~\citep{vaswani2017attention}.
A multi-layer perceptron (MLP) takes the flow time embedding as input and predicts the six modulation parameters (four scales and two biases).
The MLP is shared across all layers, different from the original DiT, and only layer-dependent biases are added to the MLP outputs.
This saves parameters without sacrificing performance. The next section describes how other inputs are conditioned.


\subsubsection{Audio Representation and Conditioning Modules}\label{sec:audio_repr_and_cond}

\textbf{Audio.}
The latent diffusion framework~\citep{rombach2021highresolution} is adopted,
where data (48kHz) is represented as compact 1D latent features of shape $T \times C$ at a much lower frame rate (25Hz) and $C=128$ extracted from a separately trained DAC-VAE (Descript Audio Codec~\citep{kumar2024high} with variational autoencoder~\citep{kingma2013auto} formulation) model.
Compared to the commonly used Encodec~\citep{Defossez2022HighFN} features (75Hz, 128-d) for 24kHz audio in audio diffusion models~\citep{shen2023naturalspeech, vyas2023audiobox},
our DAC-VAE offers a lower frame rate (75Hz$\rightarrow$25Hz), a higher audio sampling rate (24kHz$\rightarrow$48kHz), and much higher audio reconstruction quality.
Specifically, to outperform Encodec under a similar bitrate, DAC adopts the multi-scale STFT discriminators to reduce the periodicity artifacts and adds the Snake~\citep{ziyin2020neural} activation function to introduce periodic inductive biases inspired by the BigVGAN~\citep{lee2022bigvgan} architecture.
Although the code factorization technique of DAC also greatly reduces the quantization errors for much better reconstruction, discrete tokens are not necessary for diffusion-style models.
Therefore, we remove the residual vector quantizer (RVQ)~\citep{van2017neural, gray1984vector} from the DAC and trained with the variational autoencoder (VAE)~\citep{kingma2013auto} objective (which adds a KL-regularization to encourage latents to be normally distributed).
This significantly boosts the reconstruction performance especially at more compressed frame rates (25Hz).

\textbf{Video.}
Long-prompt MetaCLIP fine-tuned from MetaCLIP~\citep{xu2023demystifying} is used to extract a 1024-dimension embedding for each frame in a video.
Since the frame rate of the video might not match that of the audio, we take the nearest visual frame for each audio frame.
The resampled sequence is then projected to the DiT model dimension with a gated linear projection layer and added to the audio features frame by frame.
Adding visual and audio features frame by frame improves video-audio alignment compared to concatenating features along the time dimension,
because the former provides direct supervision of video-audio frame alignment.
We have also explored reconstruction-based features extracted from a video-autoencoder for conditioning,
which is expected to preserve more video details compared to contrastive features.
However, the results were significantly worse and also slows down training due to the large feature dimension per frame.
We concluded that the Long-prompt MetaCLIP features trained with a contrastive objective~\citep{oord2018representation, radford2021learning} encodes higher level semantic information that eases learning while keeping sufficient low-level details to capture the timing of each motion for the model to produce motion-aligned sound effects.

\textbf{Audio context.}
We follow the Voicebox~\citep{le2023voicebox} and Audiobox~\citep{vyas2023audiobox} frameworks and condition on partially masked target audio, which we coined \textit{audio context}.
This enables the model to infill (or out-fill, depending on where the mask is) audio that is coherent with the context.
Without conditioning on context, the audio would sound incoherent and change abruptly when stitching together audio segments generated independently given only the video, especially when audio contains heavy ambient sound or music.
Audio context is also represented as a DAC-VAE feature sequence, and is concatenated with the noised audio latent along the channel dimension frame by frame.
For masked frames, we replace it with zero-vectors.
To perform audio generation without any audio context, we simply input a zero-vector sequence for audio context.

\textbf{Text.}
We use text to provide additional guidance on target audio quality, sound events, and music style if music is present, which we collectively refer to as the \textit{audio caption}, details of which is described in \cref{sec:audio_cap}.
T5-base~\citep{raffel2020exploring} is used to encode an audio caption into a sequence of 768-dimensional features, where sequence length is capped at 512 tokens.
We insert a cross-attention layer right after the self-attention layer and before the feed-forward layer in each DiT transformer block for conditioning.\\

Putting it together, we denote an $N_{aud}$ frame-long sample at flow-step $t$ as $\bx_t \in \mathbb{R}^{N_{aud} \times 128}$,
and conditioning input as $\bc = \{\bc_{vid}, \bc_{ctx}, \bc_{txt}\}$,
where $\bc_{vid} \in \mathbb{R}^{N_{aud} \times 1024}$ is the resampled Long-prompt MetaCLIP features,
$\bc_{ctx} \in \mathbb{R}^{N_{aud} \times 128}$ is the masked DAC-VAE features that serves as audio context,
and $\bc_{txt} \in \mathbb{R}^{N_{txt} \times 768}$ is the T5 text feature sequence of $N_{txt}$ tokens long.
At each step, the model predicts a velocity $u(\bx_t, \bc, t; \theta) \in \mathbb{R}^{N_{aud} \times 128}$ that can be used to estimate $x_{t+\Delta t}$ during inference.
We use an additional subscript to index a subsequence when necessary. For example, $\bc_{vid,15:25}$ denotes a 10-frame segment starting on the 15th frame (0-based).


\subsubsection{Inference: One-shot Generation}
During training, each conditioning input (video, audio context, text) is dropped-out independently with some probabilities.
This enables the model to perform
\begin{enumerate*}[label=(\arabic*)]
    \item video-to-audio (V2A) generation (dropping out text and audio context),
    \item text-instructed video-to-audio (TV2A) generation (dropping out audio context),
    \item video-to-audio infilling or extension (dropping out text), and
    \item text-instructed video-to-audio infilling or extension,
\end{enumerate*}
with a single model by simply changing the conditioned inputs.

\subsubsection{Inference: Audio Extension}\label{sec:audio_ext}
Due to memory constraint and training efficiency considerations, training data is capped at a predetermined length.
To generate high quality and coherent long-form audio for videos whose lengths are beyond the cap,
we consider two algorithms: \textit{segment-level autoregressive generation} and \textit{multi-diffusion}~\citep{bar2023multidiffusion}.

Given a long video, we first split the video into \textit{overlapping} segments, and assume text captions are available for all segments.
At a high level, both algorithms run inference on individual segments first, and then consolidate the prediction.
Information can propagate across segments through overlapping frames.
Without loss of generality, we assume each segment is $n_{win}$ frames long,
and the end times of consecutive segments differ by $n_{hop}$ frames.
Two consecutive segments overlap by $n_{ctx} = n_{win} - n_{hop}$ frames.
For a video $\bc_{vid}$ of $N$ frames,
there will be $J = \lceil N / n_{hop} \rceil$ segments,
and the $j$-th segment spans $[n_{start}^{(j)}, n_{end}^{(j)}] = [max(0, (j-1)n_{hop} - n_{ctx}), min(N, jn_{hop})]$ where $j \in [J]$.
We denote the text caption for $j$-th segment as $\bc_{txt}^{(j)}$,
and the consolidated prediction at flow step $t_i$ as $\bx_{t_i}^{(j)}$,
assuming the same $T$-step flow time schedule $\{t_i\}_{i=1}^T$ are used for all segments ($t_1 = 0$, $t_T = 1$, and $t_i < t_{i+1}\;\forall i$).
Note that $\bx_0^{(j)}$ is the noise drawn from the prior $p_0$ and $\bx_1^{(j)}$ is the predicted audio for the $j$-th segment.

We introduce two extension methods below: segment-level autoregressive generation and multi-diffusion. Multi-diffusion achieves slightly better results empirically, which we use as the default method.

\textbf{Segment-level autoregressive generation.}
This algorithm emulates autoregressive generation of language models but at the segment level.
It generates one segment at a time conditioned on the information from the last $n_{ctx}$ frames of the previous segment.
Given the trajectory $\bx_{t_i}^{(j)}$ from the $j$-th segment,
we consider passing information through two routes when generating $\bx_{t_i}^{(j+1)}$: \textit{context conditioning} and \textit{trajectory regularization}.

The first route, context conditioning, is to simply update the audio context $\bc_{ctx}^{(j+1)}$ using the prediction from the previous segment.
Specifically, we set $\bc_{ctx}^{(j+1)} = [\bx_{1,-n_{ctx}:}^{(j)}; \mathbf{0}]$,
where $\bx_{1,-n_{ctx}:}^{(j)}$ denotes the last $n_{ctx}$ frames of $\bx_1^{(j)}$, and $\mathbf{0}$ is a zero matrix of shape $n_{hop} \times C$.

The second route, trajectory regularization, is to pass information through the flow trajectory $\bx_{t_i}^{(j+1)}$.
We \textit{blend} the trajectory from the previous segment with the current one at each flow step $t_i$ with a weighting function $w \in \mathbb{R}_{\ge 0}^{n_{ctx}}$,
such that the trajectory of the current segment does not diverge too much from the previous one for the overlapping segment.
Algorithmically, let $\hat{\bx}_{t_{i+1}}^{(j+1)} = \text{ODE-Solve}(u, \bx_{t_i}^{(j+1)}, \bc^{(j+1)}, t_i, t_{i+1})$ be the solution from the ODE solver at $t_{i+1}$ taking the initial condition $\bx_{t_i}^{(j+1)}$ at $t_i$ with conditioning input $\bc^{(j+1)} = \{\bc_{vid}^{(j+1)}, \bc_{ctx}^{(j+1)}, \bc_{txt}^{(j+1)}\}$ and flow model $u$.
We update the state to $\bx_{t_{i+1},0:n_{ctx}}^{(j+1)} = w \odot \hat{\bx}_{t_{i+1},0:n_{ctx}}^{(j+1)} + (1-w) \odot \bx_{t_{i+1},-n_{ctx}:}^{(j)}$ for the overlapping frames and keep other frames intact $\bx_{t_{i+1},n_{ctx}:n_{win}}^{(j+1)} = \hat{\bx}_{t_{i+1},n_{ctx}:n_{win}}^{(j+1)}$,
where $\odot$ denotes point-wise product.
In the end, we also update the audio of the overlapping frames with $\bx_{1,0:n_{ctx}}^{(j+1)}$, the consolidated prediction from the later segment.
The blending function $w$ dictates how trajectories are mixed for each frame.
To encourage a smooth transition,
we consider a linear ramping function $w_n = n / n_{ctx}$ for $n \in [n_{ctx}]$,
such that weights are higher on the previous segment for frames closer to it and vice versa.

To further boost the performance of segment-level autoregressive generation, we also explore \textit{segment-level beam search}.
At each segment, multiple candidates are generated and the resulting partial generations are ranked and pruned by a scoring model.
Those top candidates are used as prefixes to generate multiple candidates for the next segment.


\textbf{Multi-diffusion.}
Inspired by its success on leveraging a diffusion model trained on $512 \times 512$ images to generate panorama images that are 9 times wider ($512 \times 4608$) and application to video upsampling described in~\cref{sec:spatial_upsampler},
we explore multi-diffusion for audio extension, which is virtually the audio version of the panorama generation problem.
At a high level, multi-diffusion solves the ODE for each step (\eg, from $t_i$ to $t_{i+1}$) in parallel for all segments,
consolidate the prediction, and then continues with the next ODE step.
This is in contrast to segment-level autoregressive generation,
which solves the ODE for one segment completely (\ie, from $t=0$ to $t=1$, potentially with multiple steps) before solving it for the next segment.

We next describe the algorithm in detail.
To initialize, $\bx_0$ is first drawn from a prior distribution $p_0$.
At each step $t_i$, $\bx_{t_i}$ is first split into chunks $\{\bx_{t_i}^{(j)}\}_j$ based on the predefined segmentation.
$\hat{\bx}_{t_{i+1}}^{(j)}$ is computed as $\text{ODE-Solve}(u, \bx_{t_i}^{(j)}, \bc^{(j)}, t_i, t_{i+1})$ for all $j$.
Then we merge the prediction as $\bx_{t_{i+1}} = \sum_j \text{zero-pad}( m^{(j)} \odot \hat{\bx}_{t_{i+1}}^{(j)}, j)$.
The function $\text{zero-pad}(\bx^{(j)}, j)$ maps a segment of shape $n_{win} \times C$ back to a sequence of shape $N \times C$ (shape of the $\bx$),
by padding zeros based on the segment's position (\ie, the $j$-th segment spans from $n_{start}^{(j)}$ to $n_{end}^{(j)}$;
hence $n_{start}^{(j)}$ zeros are padded at the beginning and $N - n_{end}^{(j)}$ zeros at the end). 
The soft-masking function $m^{(j)} \in \mathbb{R}_{\ge> 0}^{n_{win}}$ defines how much the $j$-th segment contributes to each frame it spans.
We note that $\sum_j \text{zero-pad}( w^{(j)}, j ) = 1$, where for each frame the contribution over all segments sums to 1.
This completes one step of the flow.
The same process repeats until we reach $t_T = 1$.
Note that the weights $w$ used in segment-level autoregressive generation is applied to the overlapping frames ($n_{ctx}$ frames), 
while the soft-masking functions $\{ m^{(j)} \}_j$ are defined for all the frames each segment spans.

\cref{fig:audio_md_weights} shows two variants of soft-masking functions (zero-padded) for two segments on the right column.
The soft-masking functions $m$ are derived from window functions $\hat{m} \in \mathbb{R}_{\ge 0}^{n_{win}}$
(the left column of the same figure, also zero-padded).
Specifically, we have $\text{zero-pad}( m^{(j)}, j ) = \left( \text{zero-pad}( \hat{m}^{(j)}, j ) / \sum_{j'} \text{zero-pad}( \hat{m}^{(j')}, j' ) \right)$ given window functions $\{ \hat{m}^{(j)} \}_j$.
The uniform window function ($\hat{m}^{(j)} = \mathbf{1}$) in the top row is equivalent to the method proposed in \citet{bar2023multidiffusion}.
However, we observe that this sometimes results in abrupt transition at the boundaries (especially for smaller models),
because the prediction from two neighboring segments might still be considerably different on overlapping frames at $t=1$.
Hence the transition from taking 100\% of the first segment to 50-50 from both segments leads to discontinuity,
as shown on~\cref{fig:audio_md_weights} top row.
Using the triangular window function (\ie, Barlette window, $\hat{m}_n^{(j)} = \dfrac{2}{n_{win}-1} \left( \dfrac{n_{win}-1}{2} - \left| n - \dfrac{n_{win}-1}{2} \right| \right)$) from the bottom row results in smoother transition,
similar to what we observed in chunk-level autoregressive generation with the linear ramping function.

\begin{figure}
  \centering
  \begin{tabular}{cc}
    \includegraphics[width=0.48\textwidth]{figures/audio/md_weight_uniform_normFalse.pdf} &
    \includegraphics[width=0.48\textwidth]{figures/audio/md_weight_uniform_normTrue.pdf} \\
    \includegraphics[width=0.48\textwidth]{figures/audio/md_weight_triangle_normFalse.pdf} &
    \includegraphics[width=0.48\textwidth]{figures/audio/md_weight_triangle_normTrue.pdf}
  \end{tabular}
  \caption{\textbf{Comparing uniform weighting (top) and triangle window based weighting (bottom)} for multi-diffusion merging with an example of $l_{win}=20$ and $l_{ctx}=5$. The left column shows the unnormalized weights contributed from the first segment ([0, 15]) and the second ([10, 30]), and the right column shows the normalized weights. With triangle window, the transition in the overlapping region ([10, 15]) is much smoother.
  }
  \label{fig:audio_md_weights}
\end{figure}


\subsubsection{Model, Training, and Inference Configurations}
The DiT model has 36 layers and attention/feed-forward dimension of 4,608/18,432,
which has 13B parameters in total (excluding Long-prompt MetaCLIP, T5, and DAC-VAE).
In pre-training, videos are capped at 30 seconds (750 frames) and randomly chunked if lengths exceed.
In finetuning, we randomly sample 10s and 30s segments from the video.
Fully-sharded data parallelism is used to fit the model size.

The model is trained in two stages: \textit{pre-training} and \textit{fine-tuning}, using the same objective, but different data (described in \cref{sec:audio_pt_data} and \cref{sec:audio_ft_data}) and optimization configurations.
For pre-training, the effective batch size is 1,536 sequences, and each sequence is capped at 30 seconds (750 tokens).
The model is pre-trained for 500K updates, taking 14 days on 384 GPUs, using a constant learning rate of 1e-4 with a linear ramp-up for the first 5K steps.

For fine-tuning, the effective batch size is 256 sequences, also capped at 30 seconds.
The model fine-tunes the pre-trained checkpoint for 50K updates on 64 GPUs, which takes one day.
The learning rate linearly ramps up to 1e-4 for the first 5K steps, and then linearly decay to 1e-8 for the remaining steps.
An exponential moving average (EMA) checkpoint with a decay of 0.999 is accumulated during finetuning and used for inference.
AdamW optimizer with a weight decay of 0.1 and bf16 precision are used for both pre-training and fine-tuning.

To leverage classifier-free guidance (CFG) for inference, during training we drop conditioning inputs altogether (video, text, audio context) with a probability of 0.2.
To enable both audio generation and audio extension, the masked audio is dropped (\ie, completely masked) with a probability of 0.5, and otherwise masked between 75\% to 100\%.
To reduce the reliance on either modality, text and video input are dropped independently with a probability of 0.1 for each.

For inference, the midpoint solver with 64 steps is used.
We did not find using adaptive dopri5 solver or increasing the number of steps to boost the performance.
We use CFG with a weight of 3 on unconditional vector fields and further conduct reranking with 8 candidates per sample.
Quality score of 7.0 is used for sound effect generation, and 8.0 for joint sound effect and music generation.
For audio extension, we use dynamic guidance~\citep{wang2024analysisclassifierfreeguidanceweight} with a weight of 3 and multi-diffusion with a triangle window of $n_{win} = 40$, $n_{hop} = 30$ and $n_{ctx} = 10$ by default.

